\documentclass[10pt,a4paper]{amsart}
\usepackage[margin=19mm]{geometry}
\usepackage{amsmath,amssymb,mathtools}
\usepackage[scr=rsfs,cal=euler]{mathalfa}
\usepackage{xfrac}
\usepackage[unicode,hidelinks,bookmarks=false]{hyperref}
\newcommand{\ie}{i.e.\ }
\newcommand{\cf}{cf.\ }
\newcommand{\resp}{respectively\ }
\newcommand{\Subsection}{\subsection}
\providecommand{\nobreakdash}{\nobreak\hbox{-}}
\setlength{\emergencystretch}{2em}
\multlinegap0pt
\usepackage{amssymb}
\usepackage{tikz}
\newtheorem{lemma}{Lemma}
\newtheorem{theorem}{Theorem}
\title{Quadrangulation moves}
\author{Korablev Ph. G.}
\date{Source: arXiv:2608.01044v1 (CC BY 4.0)}
\begin{document}
\maketitle

\noindent\emph{Source and licence.} Quadrangulation moves. Version arXiv:2608.01044v1 (CC BY 4.0), distributed by arXiv under the Creative Commons Attribution 4.0 International licence, http://creativecommons.org/licenses/by/4.0/, as recorded in the arXiv licence metadata for this exact version and shown on the abstract page https://arxiv.org/abs/2608.01044v1. Full licence notice: \texttt{licenses/arxiv-2608.01044-CC-BY-4.0.txt}.\par\smallskip

\noindent\textit{Editorial scope.} Counted unit: the article's theorem Theorem:Transforms (source0.tex lines 787--790). The article's complete original proof of it is delivered with it (source0.tex lines 791--793, one \texttt{proof} environment of 3 lines). No mathematics is written, repaired or re-derived: the statement and the proof are the article's own lines, in the article's own notation, and no retained line is substituted or re-flowed.  Retained uncounted as the source-local context the proof needs: lines 451--454; lines 779--782.  Nothing was omitted from any retained span, so the package carries no omission record.  No retained line contains a citation, so the wrapper carries no bibliography and no bibliography entry.  No retained line contains a figure environment, an image-inclusion command or drawing code, so no image is delivered and the package declares no asset dependency.  Editorial formatting only, in addition: an AMS \texttt{amsart} preamble that restates the article's own declarations and macros used by the retained lines, namely the environments \texttt{lemma}, \texttt{theorem} on separate counters, together with the standard packages those lines need.  The retained environments print in this document as Lemma 1, Theorem 1 in order of appearance, whereas the article prints the same items with the numbering of its own full document.  The complete native source of the article as distributed in the arXiv e-print for this exact version is bundled unchanged as \texttt{sources/206589/source0.tex}.
\begin{lemma}
    \label{Lemma:MatchTransforms}
    Let $F$ be a surface (2-manifold), let $\mathcal{T}_1$ and $\mathcal{T}_2$ be two triangulations of $F$, and let $\mathcal{C}_1 = \beta(\mathcal{T}_1)$ and $\mathcal{C}_2 = \beta(\mathcal{T}_2)$ be quadrangulations of $F$ that match the triangulations $\mathcal{T}_1$ and $\mathcal{T}_2$ respectively. Then there exists a finite sequence of cubical Pachner moves of types $(0, 0)$, $(0, 1)$, $(1, 0)$, $(0, 2)$ and $(2, 0)$ that transforms $\mathcal{C}_1$ into $\mathcal{C}_2$.
\end{lemma}

\begin{lemma}
    \label{Lemma:Subdivide}
    Let $\mathcal{C}$ be a quadrangulation of the surface $F$, and let $\mathcal{C}'$ be the result of the subdivision of $\mathcal{C}$. Then there exists a triangulation $\mathcal{T}$ of $F$ such that $\mathcal{C}'$ can be transformed into $\beta(\mathcal{T})$ by a finite sequence of cubical Pachner moves of types $(0, 0)$, $(0, 1)$, $(1, 0)$, $(0, 2)$ and $(2, 0)$.
\end{lemma}

\begin{theorem}
    \label{Theorem:Transforms}
    Let $\mathcal{C}_1$ and $\mathcal{C}_2$ be two quadrangulations of the surface $F$. Then there is a finite sequence of subdivisions and de-subdivisions and cubical Pachner moves of types $(0, 0)$, $(0, 1)$, $(1, 0)$, $(0, 2)$ and $(2, 0)$ that transform $\mathcal{C}_1$ into $\mathcal{C}_2$.
\end{theorem}

\begin{proof}
    The theorem follows from Lemmas \ref{Lemma:MatchTransforms} and \ref{Lemma:Subdivide}. Indeed, we start from the quadrangulation $\mathcal{C}_1$, apply one subdivision, and then use the required cubical Pachner moves to obtain a quadrangulation $\mathcal{C}'$ that matches some triangulation $\mathcal{T}_1$ (by Lemma \ref{Lemma:Subdivide}). We do the same with $\mathcal{C}_2$ and obtain a quadrangulation $\mathcal{C}_2'$ that matches a triangulation $\mathcal{T}_2$. Then, by Lemma \ref{Lemma:MatchTransforms}, we can transform $\mathcal{C}'$ into $\mathcal{C}_2'$ using cubical Pachner moves, after which we apply the reverse transformations from $\mathcal{C}_2'$ back to $\mathcal{C}_2$ (the last step being a de-subdivision).
\end{proof}

\end{document}
